\scanpage{151}
Considérons maintenant des courbes elliptiques \inserted{$E$} avec ``rigidification différentielle'', définie par une base $\unsure{\varphi|\phi}$ de $t_E$, permettant donc d'identifier $t_E$ à $\mathbb C$. La donnée d'une courbe elliptique avec rigidification différentielle s'identifie donc à la donnée d'un \uline{réseau} $H\subset\mathbb C$ dans $\mathbb C$. La donnée d'une courbe elliptique avec rigidification double Torelli stricte $+$ différentielle s'identifie donc à la donnée de
\[
z_1,z_2\in\mathbb C^*\quad\text{avec}\quad
\deleted{\unsure{z_2|z}}\ \tau=z_2/z_1\in\mathcal D^+\tag*{(17)}
\]
i.e. \deleted{la $\mathbb R$ base de $\mathbb C$ \unsure{d}}\inserted{$z_1,z_2$ forment une} $\mathbb R$-base de $\mathbb C$, qui est \uline{directe} pour l'orientation canonique de $\mathbb C$. On fait opérer $SL(2,\mathbb Z)$ à gauche sur \unsure{la|leur} variété analytique complexe
\[
PD^+=\{(z_1,z_2)\mid z_1,z_2\in\mathbb C^*,\ z_2/z_1\in D^+\}\tag*{(18)}
\]
par
\[
g\,(z_1,z_2)=(dz_1+cz_2,\ bz_1+az_2)\tag*{(\deleted{18}19)}
\]
% slanted
\margin{c'est donc, via le passage $g\mapsto\tau_\infty(g^{-1})$, l'opération naturelle qui découle du Pb modulaire}

La famille modulaire bi-rigidifiée $\unsure{\mathcal X|X|\mathfrak X}_{P\unsure{\Delta|D|\mathfrak D}^+}$ peut s'écrire
% archive sheet number: 150
